% subsection: 等比型 \, $a_{n+1}=ra_n$

  \subsection{\texorpdfstring{等比型 \, $a_{n+1}=ra_n$}{等比型}}\label{節：等比型}
  \begin{bookpoint}
    \textbf{倍率を追う}\par
    一回ごとの倍率が一定なら,更新の回数だけ掛ける.公比が零の場合も,比で割らず更新の式から読む.
  \end{bookpoint}
    $r=0$なら初項は$a_1$のまま,$n\ge2$では$a_n=0$となる.
    $r\ne0$なら,$n\ge2$で掛けた回数を数えて,
    \[
      a_n=ra_{n-1}=r^2a_{n-2}=\cdots=a_1r^{n-1}.
    \]
    初項から第$n$項までに$r$を掛けるのは$n-1$回なので,指数も$n-1$になる.

    \noindent \bookblocklink{example-geometric-1}{problem}{answer}{【例題】}\par\nobreak
    \begin{enumerate}[label=(\arabic*),parsep=3pt,format=\bookitemlink{example-geometric-1}{problem}{answer}]
      \item $\displaystyle a_1=1,\, a_{n+1}=2a_n$
      \item $\displaystyle a_1=1,\, a_{n+1}=-a_n$
      \item $\displaystyle a_1=2,\, a_{n+1}=\frac{1}{2}a_n$
    \end{enumerate}

    \noindent\bookblocklink{example-geometric-1}{hint}{problem}{【方針】}\par\nobreak
    \begin{enumerate}[label=(\arabic*),parsep=3pt,format=\bookitemlink{example-geometric-1}{hint}{problem}]
      \item 等比型の\bookterm{general}{一般項}
            $\displaystyle a_n=a_1r^{n-1}$
            を用います.
      \item 公比は$r=-1$です.
            符号が交互に変化することに注意します.
      \item 初項と公比を\bookterm{general}{一般項}に代入します.
    \end{enumerate}

    \noindent\bookblocklink{example-geometric-1}{answer}{problem}{【解答】}\par\nobreak
    \begin{enumerate}[label=(\arabic*),parsep=3pt,format=\bookitemlink{example-geometric-1}{answer}{problem}]
      \item $\displaystyle a_n=2^{n-1}$
      \item $\displaystyle a_n=(-1)^{n-1}$
      \item $\displaystyle a_n=2\left(\frac12\right)^{n-1}$
    \end{enumerate}

    \noindent\bookblocklink{example-geometric-1}{solution}{problem}{【解法】}\par\nobreak
    (1)は$a_1=1,\ r=2$なので$a_n=2^{n-1}$となる.$2^0=1$で初項が一致し,$2^n=2\cdot2^{n-1}$で更新式も満たす.
    (2)のように公比が負でも,掛ける回数は変わらない.
